\documentclass[11pt,a4paper]{article}
\usepackage[margin=22mm]{geometry}
\usepackage{amsmath,amssymb,amsthm,booktabs,microtype,needspace}
\usepackage[T1]{fontenc}
\usepackage{lmodern}
\usepackage[colorlinks=true,linkcolor=blue,urlcolor=blue]{hyperref}
\hypersetup{pdftitle={Subset sums: core cutoffs and the audited bound 92},pdfauthor={Sum-of-Subset research project}}
\newcommand{\FS}{\operatorname{FS}}
\newcommand{\Z}{\mathbb Z}
\newcommand{\T}{\mathcal T}
\newtheorem{theorem}{Theorem}[section]
\newtheorem{lemma}[theorem]{Lemma}
% Exact-run statistics are filled from the verified result before publication.
\newcommand{\RootGroups}{14,700}
\newcommand{\CoarseGroups}{14,518}
\newcommand{\NestedGroups}{182}
\newcommand{\IntervalChecks}{14,761}
\newcommand{\ExtremumPoints}{35,013}
\newcommand{\TreeRoots}{532}
\newcommand{\BinaryNodes}{1,706}
\newcommand{\TreeLeaves}{2,238}
\newcommand{\MaxDepth}{13}
\newcommand{\FinalSystems}{128,106}
\newcommand{\Vectors}{2,450}
\newcommand{\MinGap}{1/936}
\newcommand{\NegativeControls}{31}
\title{Subset Sums in a Fourfold Sumset\\\large Core cutoffs and the audited bound $N\le92$}
\author{Sum-of-Subset research project}
\date{16 October 2026}
\begin{document}
\maketitle
\begin{abstract}
We continue the independently audited chain from $N\le93$ to $N\le92$.
At cardinality 93, the paired bootstrap only guarantees positive-side
size 13; we first eliminate this boundary, and use the previous exact
prefix inverse already in the lattice gate. For the covering argument,
positive low holes are separated into three bands before mass transport.
Their endpoints are 2, 4 and 14, with respective covering caps 262, 254
and 250. This preserves the newly possible three-empty-unit family for
the first band while avoiding the unnecessary transport to 14 there.
A unit-coefficient propagation proves outer integrality in 2,987 of
the 2,992 three-empty modes. A new core cutoff retains the second-largest
element on each positive side and extends the two-tail bound from 25
to 36 cutoff pairs. Core mass lower bounds and a dense-span
lower bound remove incompatible combinations of small span and many
holes. A new persistent-hole argument bounds the defect directly,
raising a critical pair of floors from 97 and 126 to 112 and 144.
High-hole groups retain one common hole variable in every coupled row.
Fresh integer certificates close the complete endpoint. The stronger
archived bound 84 still awaits a complete independent replay; the
present result advances the audited chain.
\end{abstract}

\section{Result and inherited lemmas}
For a finite real set $A$, write
$\FS(A)=\{\sum_{x\in X}x:X\subseteq A\}$, including zero, and
$4A=A+A+A+A$, allowing repetitions. Let $N$ denote the maximum feasible
size under $\FS(A)\subseteq4A$.
\begin{theorem}\label{main}
Every feasible set has at most 92 elements. Hence $71\le N\le92$.
\end{theorem}
The unchanged 71-element construction is
$([-31,31]\cap\Z)\cup\{\pm93,\pm155,\pm217,\pm713\}$.
The inherited audited bound is $N\le93$ from \cite{previous}.
It therefore suffices to exclude exactly 93 elements. Such a set contains
zero: otherwise adjoining zero preserves feasibility and produces 94
elements. No inheritance of feasibility by arbitrary deletion is used.

We use the corrected integralization, peeling, simultaneous inverse,
effective-hole, mass, span and covering lemmas of \cite{foundation},
the sharp four-extreme and joint-covering results of \cite{overlap},
and the non-progression and prefix results of \cite{sharp,previous}.
For a positive $r$-set $X$, put $T_r=r(r+1)/2$ and
\[
 \mu(X)=|\FS(X)|-1-T_r,\qquad
 \delta_0(X)=|\{x\in X:2x\notin R(X\cup\{0\})\}|,\qquad
 \rho(X)=\mu(X)+2-2\delta_0(X),
\]
where $R$ uses distinct-element two-sums. The inherited general floors
are $\rho\ge-8$, and $\rho\ge0$ for $r\ge14$.
If the subset-sum support is not a progression, $\mu\ge r-3$,
$\rho\ge r-14$ for $r\ge6$, and $\rho\ge r-5$ for $r\ge17$.
The last threshold is not extended to 16 here.

Let $F(r,u)$ denote the exact integer prefix floor proved in
\cite{previous}: for $r\ge17$, non-progression support and $\mu\ge u$
imply $\rho\ge F(r,u)$. Its complete formula and proof are inherited;
\path{verification/nonap_inverse.py} is unchanged. The 2,723 base cases
and 5,409 integer inverse checks are rerun. The new hole-dependent
estimate below strengthens $F$ without changing it.

\section{A stricter order for the size-93 lattice gate}
Remove four original extremes on either side to obtain $D$ and peel
twice to $D\supset D_1\supset E=P\cup\{0\}\cup(-Q)$. Write
\begin{gather*}
 d=85,\quad e=77,\quad a+b=m=76,\quad I=(a-38)^2,\\
 \Lambda=1+T_a+T_b=1483+I,\qquad \Omega=484-I.
\end{gather*}
The inherited identity is
\[
 \rho(P)+\rho(Q)+2\vartheta+c+\T=\Omega,
 \qquad c=2|2D|+18d+85-|\FS(D)|,
\]
where $\vartheta,c,\T$ are nonnegative and
$\T=\kappa+\ell_1+\ell_2+g$.
The deficient-sign test $1+T_{77}>24\cdot93+77$ holds.
The paired bootstrap gives $q^2\le2536$ with even $q$, so $q\le50$
and $a,b\ge13$. If one side is 13, the other is 63 and
$\Omega=-141<-8$. This contradicts the identity using the general
$-8$ floor on the small side and the zero floor on the large side.
Now both sides are at least 14, so $\Omega\ge0$, leaving
$16\le a\le60$. Across all 45 orientations,
\[
 T_{\min(a,b)}-2-\Omega\ge8>0.
\]
The simultaneous inverse supplies a common core lattice. Define
$z=h_P+h_Q-\kappa$, so $0\le z\le\min(h_P,h_Q)\le\kappa$, with
$h_X<|\FS(X)|-1$.

Put $\rho=\rho(P)+\rho(Q)$, $\mu=\mu(P)+\mu(Q)$ and
$C_0=2m+4e+14-\Lambda=-1009-I$. The two inner comparisons are
\[
 |\FS(D_1)|-2|\FS(E)|\le
 \min\{C_0+\Omega+\rho-\mu+2z,\ C_0+3\Omega-\rho-\mu\}.
\]
Since the maximum is defective, $\rho\le\mu$. Thus $z\le14$ gives
an upper bound $-497-2I<0$.
For $z\ge15$, both supports are non-progressions. First use
$\mu\ge70$ and the applicable side floors $r-5$ or, at size 16,
$r-14$. If that comparison is not yet negative, the parameters are
$30\le a,b\le46$. Use the effective-hole witness at $\beta=7$ on
each side to obtain $u_r=r+C_7(r)+1$, and substitute the stronger
\emph{combined} estimate
\[
 \rho+\mu\ge F(a,u_a)+F(b,u_b)+u_a+u_b.
\]
All witness tests are exact. The resulting high-hole comparison is at
most $-17$. At the balanced parameter, the earlier floor alone gives
$443-(66+310)=67$, which fails; the prefix floor gives
$443-(206+310)=-73$. This use of $F$ precedes the outer lattice argument
and is valid for positive real sets, so it does not assume the conclusion.
The outer comparison is
\[
 \Lambda-(4d-10+\Omega)=669+2I>0.
\]
An off-lattice element would give disjoint support translates and
violate the corresponding strict factor-two comparison. Hence $D_1$
and $D$ lie on the core lattice. Only after this step do we use hole
propagation and the sharp four-extreme lemma:
\begin{equation}\label{budget}
 w\le4z,\qquad\T\ge\kappa+160-3z,\qquad
 \T\ge\max(54,z,160-2z).
\end{equation}
At the two size-16 boundary orientations $\Omega=0$, so no state
survives \eqref{budget}. The prefix inverse is therefore never called
outside its proved range in the subsequent covering models.

\section{A direct defect bound from persistent holes}
\begin{theorem}\label{holerho}
Let $X$ be a positive integer $r$-set with hull-hole count
$2\beta<h_X<|\FS(X)|-1$. Suppose
\[
 \beta\ge2,\qquad r\ge L\ge\max(17,2\beta+5).
\]
Then $\rho(X)\ge B(r,\beta,L)$, where
\begin{equation}\label{holerhofloor}
 B(r,\beta,L)=\min\{(\beta-1)(r-L+1),\ T_L-r-1,
                   \ T_L+L-6\beta-r-7\}.
\end{equation}
\end{theorem}
\begin{proof}
Assume $\rho(X)$ is below all three terms. The hole condition makes $X$
primitive: a gcd at least two gives at least $|\FS(X)|-1$ hull holes.
Its support is not a progression, since a primitive progression support
containing zero would be a full interval.
Write $X_i$ for increasing prefixes and
$y_i=|\FS(X_i)|-|\FS(X_{i-1})|-i$. We use
$\rho_i\ge\rho_{i-1}+y_i-2$ and the inherited fact that, for $i\ge6$,
$y_i=0$ implies a progression support.

If the last progression prefix had size $t\ge L$, its next translate
would be disjoint, so $y_{t+1}\ge T_t-t$. Its predecessor has
$\rho_t\ge0$; every later prefix is non-progression and has $y_i\ge1$.
Consequently
\[
 \rho_r\ge T_t-t-2-(r-t-1)=T_t-r-1\ge T_L-r-1,
\]
a contradiction. All prefixes of size at least $L$ are therefore
non-progression. If every $y_i>\beta$ for $L\le i\le r$, telescoping
from $\rho_{L-1}\ge0$ gives
$\rho_r\ge(\beta-1)(r-L+1)$, another contradiction.
Choose $i\ge L$ with $1\le y_i\le\beta$.

The small-increment inverse writes $X_i=tY$, where $Y$ is a primitive
deleted interval with at most $\beta$ deletions. The scale $t$ is an
integer because $X_i$ is integral and $\gcd(Y)=1$.
If $t>1$, final primitivity forces a first later element off $t\Z$,
at index $k$. Its support translate is disjoint, and the same estimate
with $\rho_{k-1}\ge0$ and later $y_j\ge1$ gives
$\rho_r\ge T_{k-1}-r-1\ge T_L-r-1$. Hence $t=1$.
Since $i\ge2\beta+5$, the inherited counted-endpoint lemma applies:
\[
 \FS(X_i)=[0,S]_{\Z}\setminus(H\cup(S-H)),\quad
 H\subseteq[1,2\beta],\quad |H|\le\beta.
\]

Suppose this description holds at size $v\ge i$, with sum $S$.
Let $x$ be the next element. If $x\ge S-4\beta$, its translate supplies
at least $S-6\beta$ new points above $S$: if $x\le S$, at least
$x-2\beta$ of the last $x$ positions are present; if $x>S$, the whole
translate is new. Thus
$y_{v+1}\ge S-6\beta-v-1$. The sharp non-progression floor
$\rho_v\ge v-5$ and the remaining increments $y_j\ge1$ give
\[
 \rho_r\ge S+v-6\beta-r-7
 \ge T_L+L-6\beta-r-7,
\]
contradicting the third term. Therefore $x<S-4\beta$.
Also $x>2\beta$, since it exceeds the maximum of a prefix of size
at least $2\beta+5$. The two central support intervals overlap or abut;
the lower endpoint holes stay $H$, and complement symmetry gives the
upper holes. Induction preserves this description to $X_r$, which has
at most $2\beta$ holes. This contradicts the hypothesis.
\end{proof}
Define $\mathcal R(r,z)$ as the maximum of zero and the accepted
$B(r,\beta,L)$ for $2\beta<z$, searching $\beta\le20$ and all admissible
$L\le r$. Set
\[
 W(r,u,z)=\max\{F(r,u),\mathcal R(r,z)\}.
\]
For example, when $z\ge19$, the choices
$(r,\beta,L)=(36,8,21)$ and $(40,9,23)$ give floors 112 and 144.
At excess floors $u_p=148$, $u_q=184$, the old prefix floors were
97 and 126. The direct hole argument raises their sum by 33.
The new helper \path{verification/nonap_hole_rho.py} checks all
10,869 searched parameter rows and enforces the strict hole threshold.
No numerical optimization is used to justify this theorem.

\section{Bandwise mass transport}
The exact effective-hole witness is recalled to specify the finite input.
For integers $r,\beta\ge1$, $C\ge0$, let $K_0\ge3$ be least with
$(K_0+1)(K_0-2)/2>C+1$, and put
\begin{gather*}
 D_0=C+K_0,\quad E_0=D_0+1,\\
 M_0=\max\left(2\beta,1+\left\lfloor
 \frac{1+\sqrt{9+8(E_0+6\beta)}}2\right\rfloor\right),\quad
 J=\max(K_0+1,2\beta+5,M_0).
\end{gather*}
Accept a witness only if
\[
 T_{M_0}-6\beta>M_0+1+E_0,\quad
 \beta(r-J+1)>D_0,\quad T_{J-1}-1>D_0.
\]
Then $2\beta<h_X<|\FS(X)|-1$ implies $\mu(X)\ge r+C+1$.
The catalogue searches $17\le r\le59$, $1\le\beta\le20$,
$0\le C\le600$. These limits need not be optimal; every used witness
is checked independently of the search.
Let $U(r,z)$ be the maximum of $r-3$ and the accepted $r+C+1$ with
$2\beta<z$. For $z=1,2$ use $U=r-3$.
For positive holes, \eqref{budget} gives
\begin{equation}\label{cap}
 c\le484-I-W(a,U(a,z),z)-W(b,U(b,z),z)-\max(54,z,160-2z).
\end{equation}
For zero holes use zero side penalties, obtaining $c\le324-I$.
All surviving integer states are enumerated within $0\le z\le\Omega$.
There are 443 low-hole states and 3,602 high-hole states.

\begin{lemma}[Three bands]\label{bands}
For positive low holes, the following bounds hold before transport.
\begin{center}
\begin{tabular}{lrrrr}\toprule
Band&Hole range&Combined $\rho$ floor&$c+I$ cap&Endpoint\\\midrule
$P$&1--2&66&262&2\\
$Q$&3--4&78&254&4\\
$R$&5--14&102&250&14\\\bottomrule
\end{tabular}
\end{center}
\end{lemma}
\begin{proof}
The first floor is $(a-5)+(b-5)=66$. For the second, the $\beta=1$
witnesses and $F$ give a combined floor at least 78 for every
$17\le a\le59$; the $\beta=2$ witnesses give at least 102 for the third.
All 86 witness pairs and exact inverse parameters are recorded.
Since $z\le14$, the last term in \eqref{cap} is $160-2z$.
Thus $c+I\le324-\rho+2z$, giving exactly the table.
\end{proof}
In particular, $262\ge3d=255$ but $254,250<255$.
The three-empty family is possible only in band $P$ and must be retained.
A coarse use of $\rho\ge66$ for all low holes would give cap 286.
We do not claim fewer total roots from banding: the stronger models
have 52,660 positive integer roots versus 47,212 for that weaker
one-band relaxation. The advantage is the smaller transport error and
the stronger valid budgets, not that comparison of root counts.

\begin{lemma}[Endpoint transport]\label{transport}
Fix a band with upper endpoint $Z$ and impose its constant cap first.
Every actual model in the band maps to its endpoint model by
\[
 z\leftarrow Z,\qquad S_p\leftarrow S_p+Z-z,\qquad
 S_q\leftarrow S_q+Z-z,
\]
leaving other coordinates fixed.
\end{lemma}
\begin{proof}
Write a model row as $v\cdot x\le b$ and $t=Z-z\ge0$.
Its left side changes by
$t(v_{S_p}+v_{S_q}+v_z)$. Every retained row except the upper-hole
row has this coefficient nonpositive, as checked symbolically by the
verifier. The upper-hole row holds by equality at $Z$; all coordinates
remain nonnegative. The cap is already independent of $z$, and its
right side is not recomputed after transport. The lower-hole row also
weakens. Therefore all rows hold at the endpoint. The inherited endpoint
binary alternatives are selected afresh using the enlarged mass endpoint.
\end{proof}
The largest transport increments are now 1, 1 and 9 on the three bands.
Sending all positive low holes to 14 would permit increments up to 13,
including the delicate newly restored three-empty family.

Let $j$ count empty units and put $\delta=c-85j$.
The low-hole root ranges, including all retained modes, are
\begin{center}\small
\begin{tabular}{lrr@{\qquad}lrr}\toprule
Regime&$a$ range&Integer roots&Regime&$a$ range&Integer roots\\\midrule
$Z_0$&20--56&148&$P_0$&22--54&132\\
$Z_1$&23--53&2,170&$P_1$&25--51&1,890\\
$Z_2$&26--50&14,450&$P_2$&29--47&10,982\\
$Z_3$&30--46&50,864&$P_3$&36--40&14,960\\
$Q_0$&23--53&124&$R_0$&23--53&124\\
$Q_1$&25--51&1,890&$R_1$&26--50&1,750\\
$Q_2$&29--47&10,982&$R_2$&30--46&9,826\\\bottomrule
\end{tabular}
\end{center}
Here $Z$ denotes zero holes. This is a containing universe: every
surviving actual scalar state lies in the listed range, although some
listed parameters may admit a stronger scalar exclusion.

\section{Recovering integrality and extending the core cutoff}
\begin{lemma}[Integral-center propagation]\label{integrality}
Suppose the group residual cap is below $2d-1=169$.
Start from the two leading centers, all selected nonempty ordinary
centers, and the actual member of any split pair. Take differences of
known integral forms and repeatedly solve a form with one unknown
coordinate of coefficient $+1$ or $-1$. Every recovered coordinate is
integral on the core lattice.
\end{lemma}
\begin{proof}
Both leading losses are at most $\delta<169\le|2D|$, so each leading
translate contributes. Because $D$ and $2D$ are integral, contribution
to the integral target forces its center to be integral. The same argument
applies to every selected nonempty onefold translate and actual split-pair
member. At most three missing units leave at least one own-center slot
on each side, so both leading losses are included in their actual block
deficits and are bounded by $\delta$.
Differences of integral forms remain integral. In a form
$\pm x+\sum c_i x_i\in\Z$ with integral $c_i$ and already integral
$x_i$, subtraction proves $x\in\Z$; induction proves every step.
Division by a coefficient other than $+1$ or $-1$ is not allowed.
\end{proof}
The checker begins with
$c_p=p_1-p_2-p_3-p_4$, $c_q=q_1-q_2-q_3-q_4$.
For every selected center $v$ it also forms $v-c_p$ or $v-c_q$.
For example, with branch A and missing labels $p0,p1,q5$, the retained
$p2,p3$ centers recover $p_3,p_4$. The center
$p8=p_2-p_3-p_4$ then recovers $p_2$, and the leading identity recovers
$p_1$. The other side follows from its retained own centers.
All eight coordinates are integral, so every original outer gap is at
least one; in branch B its strictly positive integral sum is at least one.

Complete propagation succeeds in 2,987 of the 2,992 three-empty modes.
The remaining five retain weak outer gaps. They all lie in branch B:
\begin{center}\small
\begin{tabular}{lll}\toprule
Missing labels&Selected second-pair member&Unrecovered coordinates\\\midrule
$p0,p1,p8$&$p$&$p_1,p_2$\\
$p1,p2,q6$&$p$&$p_2,p_3$\\
$p1,p3,q7$&$p$&$p_2,p_4$\\
$p6,q1,q2$&$q$&$q_2,q_3$\\
$p7,q1,q3$&$q$&$q_2,q_4$\\\bottomrule
\end{tabular}
\end{center}
This list describes what the stated propagation rule proves, not a claim
that the five exceptional modes must have nonintegral outer elements.
The verifier recomputes all modes and each unit-coefficient derivation;
no integrality declaration is read from a certificate.

The decisive geometric refinement is to retain $C_p$, the second-largest
element of $P$, and $C_q$, the second-largest element of $Q$.
On side $p$, with $x_0=R_p$ and $x_4=M_p$ as below, actual sets satisfy
\begin{equation}\label{coremass}
 C_p\ge a-1,\qquad C_p+1\le R_p,\qquad
 C_p+R_p+x_1+x_2+x_3+M_p-S_p\le-T_{a-2}.
\end{equation}
The last row counts the other $a-2$ distinct positive integers at their
minimum possible sum. It remains valid after mass transport, which only
increases $S_p$. Reflect for side $q$.
We also use the upper mass constraint
\begin{equation}\label{coreupper}
 S_p\le x_1+x_2+x_3+M_p+R_p+(a-1)C_p-T_{a-2}.
\end{equation}
Indeed, the other $a-1$ core elements are distinct integers at most $C_p$.
To keep a constant matrix over an $a$-interval, the verifier replaces
$a-1$ in its coefficient by the group's largest possible $a-1$; this
weakens the upper bound since $C_p\ge0$. For a positive low band
$[z_0,Z]$, use the invariant row obtained by subtracting $z$ on the left
and $z_0$ on the right. It holds before transport because $z\ge z_0$,
and $S_p-z$ is unchanged by transport. At the endpoint it allows only
the justified extra mass $Z-z_0$. High groups retain the original row
and their exact fixed side size.


\begin{lemma}[Independent core cutoffs]\label{cutoff}
Put $x_{-1}=C_p$, $y_{-1}=C_q$. For every $-1\le i,h\le4$,
\begin{equation}\label{tails}
 |2D|\le M_p+M_q+x_i+y_h+1+T_{4-i}+T_{4-h}.
\end{equation}
\end{lemma}
\begin{proof}
The central integer interval $[-M_q-y_h,M_p+x_i]$ has
$M_p+M_q+x_i+y_h+1$ points. A sum above its upper endpoint must use
two positive elements strictly above $x_i$: otherwise one summand is
at most $x_i$ and the other at most $M_p$. There are $4-i$ such elements,
so there are at most $T_{4-i}$ unordered pairs allowing repetition.
Likewise at most $T_{4-h}$ distinct sums lie below the lower endpoint.
Adding these three bounds proves the claim, independently for each cutoff.
\end{proof}
The old model used only $0\le i,h\le4$. The added positions give
36 pairs, including 11 new rows, and two new coordinates constrained
by \eqref{coremass}. For illustration, the previously surviving
size-93 relaxation at $a=37$, $b=39$ had
\[
 R_p=73,\ R_q=48,\ M_p=200,\ M_q=110,\ S_p=1165,
 \quad (x_1,x_2,x_3)=(74,75,77).
\]
Its mass row forces $C_p\le36$. The new cutoff $(i,h)=(-1,0)$ therefore
gives $|2D|\le200+110+36+48+1+15+10=420$, instead of the relaxed value
447. This is an excluded relaxation point, not a construction.
The retained inner-loss disjunctions now extend to the six largest
magnitudes $M_p,x_3,x_2,x_1,R_p,C_p$, and similarly on side $q$.
If fewer than $i$ points are lost at the lower endpoint, a contributing
center is at least one above the $i$th negative magnitude; the upper
version uses the $i$th positive magnitude and the mass endpoint.
Thus the new two positions use the same counting proof as the first four.

For a finite sanity check, take each six-element subset of
$\{1,2,3,40,41,42,43,100\}$. For all 784 ordered pairs, independently
compute the actual twofold sumset and check all 36 bounds and the counts
outside their central intervals: 28,224 checks, with 159 strict
improvements over the old cutoff minimum. The arbitrary-set
statement is proved above; these finite checks do not replace that proof.

\section{Dense spans and the high-hole models}
\begin{lemma}[Holes require span]\label{density}
Let $X$ be a positive integer $r$-set, $r\ge4$, with maximum $R$ and
$h$ holes in its subset-sum hull. If $0\le z\le h$, then
\begin{equation}\label{densitybound}
 2R\ge2r+\min(z,r-3).
\end{equation}
\end{lemma}
\begin{proof}
If $2R\ge3r-3$, the conclusion is immediate. Otherwise integrality gives
$2R\le3r-4$. Put $q=R-r$; there are exactly $q$ deletions from $[1,R]$.
Lev's central-interval theorem \cite[Theorem 1]{Lev1998}, already used
in the inherited hole lemmas, applies under $R\le3r/2-2$ and gives
$[2q+1,\sum X-(2q+1)]\subseteq\FS(X)$.
Each missing sum in $[1,2q]$ is a deleted element, so there are at most
$q$ such sums. Complement symmetry gives $h\le2q$; the two end regions
are disjoint since $\sum X\ge T_r>4q$. Hence $z\le2(R-r)$, proving
the other sufficient inequality and therefore \eqref{densitybound}.
\end{proof}
For the three low bands, their lower endpoints satisfy $z_0\le r-3$.
Integrality therefore gives the transport-invariant rows
$R\ge r+\lceil z_0/2\rceil$, adding respectively 1, 2 and 3 to the
former $R\ge r$ floor. For a high group with hole interval $[L,H]$, keep
the endpoint floor $R\ge r+\lceil\min(L,r-3)/2\rceil$ and the chord
of the concave function $f(z)=\min(z,r-3)$. With
$l=f(L)$, $u=f(H)$ and $H>L$, this is the integer-coefficient row
\[
 2(H-L)(R-r)\ge l(H-L)+(u-l)(z-L).
\]
When needed, an exact binary alternative keeps either
$z\le r-3$ and $2R\ge2r+z$, or
$z\ge r-3$ and $2R\ge3r-3$. Both children are checked.
The chord uses the same $z$ as every covering row, not a separately
chosen extremum. The group audit freezes the target chord endpoints
when comparing each source state to its containing model.

The checker directly enumerates 16,593 dense deleted intervals of sizes
6--14, verifies the hole bound, and checks 14,491 integer chord points.
A non-dense even progression is a negative control against dropping the
density condition. These checks supplement the general proof above.

For $z\ge15$, the scalar cap is at most $174<3d$.
The first pair's actual contributor fixes the direction in the inherited
active-center test; reflection cannot be chosen to improve the test.
There remain 1,330 states with $j=0$, 1,535 with $j=1$, and 24 with
$j=2$. Group by $(a,j,U(a,z),U(b,z),v_p,v_q)$, where the two $W$ floors
$v_p,v_q$ are constant, and retain a common real variable
$z\in[z_{\min},z_{\max}]$. There are respectively 17, 45 and 5 groups.
Write $E_*=484-I-v_p-v_q$. In addition to
$\delta\le C_*-85j$, retain
\[
 \delta+z\le E_*-85j,\quad
 \delta\le E_*-54-85j,\quad
 \delta-2z\le E_*-160-85j.
\]
The same $z$ occurs in all mass, span and cardinality rows. Missing
integer values and real intermediate values may be included in the
relaxation, but the cap and hole count may not be independently maximized.
The group audit also verifies the common $W$ floors using the prefix
floor and the hole witness at the group's minimum $z$. It checks 947,062 row comparisons over all 2,889
retained states, using the target group's fixed row policy.
The high-hole root count becomes
$17\cdot4+45\cdot70+5\cdot578=6,108$.

We retain the twenty-center, eighteen-unit grammar from
\cite{overlap,previous}; the numbers of modes for $j=0,1,2,3$ are
$4,70,578,2992$. Every choice of missing units is enumerated.
For a potentially split pair both actual-contributor choices are kept;
positivity is used only as a necessary condition for contribution.
Write $x_0=R_p<x_1<x_2<x_3<x_4=M_p$, and similarly $y_i$, and
$K=|2D|$. The mass rows on side $p$ are
\begin{align*}
2M_p+6R_p-2S_p+2z&\le-2T_a-12-2u_p,\\
M_p+4R_p-S_p&\le-T_{a-1}-6,\\
x_1+x_2+x_3+M_p-S_p+z&\le-T_a-u_p,\\
R_p+x_1+x_2+x_3+M_p-S_p&\le-T_{a-1}.
\end{align*}
Use $u_p=0$ for $Z$, $u_p=a-3$ in each positive low band, and the
constant $u_p=U(a,z)$ in a high group; reflection supplies side $q$.
The gap rows include $R_p\ge a$, $M_p-R_p\ge4$ and
$x_i-x_{i-1}\ge1$. The numerical span and cardinality rows are
\[
 R_p+R_q+\delta-4z\le404-I-85j,\qquad
 S_p+S_q+\delta-2K-4z\le1614-85j.
\]
Low-band caps are $324-I-85j$ for $Z$ and the corresponding table cap
minus $I+85j$ otherwise. High groups use the coupled rows above.

The group-wide residual cap $C_g$ determines all conditional rows:
\begin{center}\small
\begin{tabular}{ll}\toprule
Rows&Required condition\\\midrule
Ordinary quantiles&$C_g<r_{\min}+5$\\
Core quantiles&$C_g<81$\\
Basic leading quantiles&$C_g<169$\\
Nested leading quantiles&$C_g<161$\\
Strict outer gaps&Inherited $j\le2$ condition, or Lemma~\ref{integrality}\\\bottomrule
\end{tabular}
\end{center}
The inherited strict-gap condition is $j\le1$, or $j=2$ with
$C_g<169$. Here $r_{\min}$ is the smallest side size in the entire group. Use the
strongest justified individual type. For a positive center $v$, global
quantiles are $v+l+M_p\ge85$ and $v-u-M_q\le S_p-84$.
Core versions replace $(M_p,M_q,85,84)$ by $(R_p,R_q,81,80)$.
Ordinary versions are $v+l\ge b+5$ and $v-u\le S_p-a-4$.
Individual loss caps are 84 globally and 80 for core rows.
In particular, $Z_3$ has residual cap 69 and uses core rows;
$P_3$ has cap 7 and uses ordinary rows. For both three-empty families,
Lemma~\ref{integrality} licenses strict outer gaps in the 2,987 recovered
modes; the five unrecovered modes retain weak gaps.
High one-empty residual caps can exceed 81; no universal core shortcut
is taken. High two-empty residual caps are at most 4 and hence justify
strict outer gaps, without automatically justifying ordinary quantiles.

As in \cite{overlap}, each block deficit $O$ is charged once. For a
block of $k$ ordinary centers,
\[
 O\ge77k-(R_p+R_q+v_{\max}-v_{\min}+1)+\max_v(l_v+u_v).
\]
For a leading slot and a nonempty subset $S$ of its retained own-center
set $V$, use
\[
 O\ge77(|S|+2)-2-p_{\max}-R_p-2R_q
 +\max\{L_p,l_v+u_v:v\in S\}
 +\sum_{v\in V\setminus S}(l_v+u_v).
\]
Also impose the individual-loss sum bound; in an own block the leading
loss is added to the ordinary hull bound. These are alternative lower
bounds on the same deficit, not independent gains to be added. Sum only
disjoint blocks and ungrouped losses, to at most $\delta$.
Leading quantiles use $2d-1=169$ and $2e+7+2i=161+2i$ under the cap
conditions above. The coarse model substitutes
$K\le M_p+M_q+R_p+R_q+21$. The nested model keeps five telescoping envelope levels and all
36 two-tail inequalities \eqref{tails}, together with \eqref{coremass}.
The coarse model eliminates $C_p,C_q$ and their rows; the extra core
information is introduced only when the coarse contradiction fails.
All coordinates are nonnegative reals. The four inherited binary
alternatives remain available: core quantiles or loss at least 81;
ordinary quantiles or an outside endpoint; an inner loss count or its
cutoff; and both or a sole contributor in the split pair. Every binary
node must certify both children. Multiple large-loss centers are allowed. High groups additionally use
the density alternative proved above as an optional discovery refinement.
The final certificate pool uses neither that split nor the upper-mass
rows \eqref{coreupper}; those optional rows are not dependencies of the
final contradiction.

\section{Exact coverage and reproducibility}
The checker rebuilds every root and row before reading certificate
assignments. There are $120292+6108=126400$ integer root systems.
The verified proof statistics are
\begin{center}\small
\begin{tabular}{lr}\toprule
Quantity&Count\\\midrule
Structural groups&\RootGroups\\
Coarse / nested groups&\CoarseGroups\ / \NestedGroups\\
Symbolic interval checks&\IntervalChecks\\
Integer extremum evaluations&\ExtremumPoints\\
Integer roots requiring trees&\TreeRoots\\
Binary internal nodes / leaves&\BinaryNodes\ / \TreeLeaves\\
Maximum depth&\MaxDepth\\
Final integer leaf systems&\FinalSystems\\
Distinct integer vectors&\Vectors\\\bottomrule
\end{tabular}
\end{center}
For $Ax\le b$, $x\ge0$, a nonnegative integer multiplier $w$ with
$w^TA\ge0$ and $w^Tb<0$ is an exact contradiction. The smallest
normalized gap $-w^Tb/\sum w_i$ is $\MinGap$.
Of the distinct vectors, 1,354 use a new core cutoff, 434 use a density
chord, and 113 use a sixth-position inner-loss row; these counts may
overlap and are recorded in \path{method_usage.json}.
Symbolic right sides are constructed as quadratics directly; their
integer extrema are checked at endpoints and at the integers around
the vertex. Each integer parameter is assigned exactly once. Tree leaves
are checked separately, with both children required and repeated or
unproved splits rejected. Unused multipliers are rejected as well.
The \NegativeControls\ negative controls include omission of the newly
required $P_3$ family, omission of high two-empty modes, invalid row
thresholds, unproved or universal outer integrality, an unjustified density bound, arbitrary reflection, incompatible hole/cap maxima and an
attempt to exclude the next cardinality 92.

\subsection*{Reproduction and scope}
Only Python 3.10+ and its standard library are needed for checking:
\begin{verbatim}
python -I -S -B verification/upper92.py --output .work/upper92.json
python -I -S -B verification/upper92.py --exhaustive
\end{verbatim}
The certificate data are \path{data/upper92_certificates.json}.
Exhaustive mode checks every integer $a$, every numeric row against
its symbolic expression, and every multiplier pointwise. In a high group,
$z$ remains a coordinate; each exact contradiction covers its entire
real interval simultaneously. The optional generator
\path{experiments/generate_upper92_certificates.py} uses SciPy only
for discovery and accepts certificates only after exact checking.
Standard and optimized Python results are byte-identical. Exhaustive
and interval results agree except for the mode label; regeneration is
byte-identical, and the upper-93 baseline result is unchanged.
The full result, witness data, nested-group index, generation logs,
hole-defect witnesses, input hashes and run receipt are saved in
\path{results/2026-10-16/upper92/}.
\path{verify_upper92_windows.bat} runs all ten finite continuations.
These computations verify the declared finite reductions, conditional
on the general mathematical lemmas proved here and in the foundation;
they are not a proof-assistant formalization. The audited conclusion is
$N\le92$. Cardinality 92 requires a new lattice and coverage analysis.

\begin{thebibliography}{9}\raggedright
\bibitem{previous} Sum-of-Subset research project,
\emph{An integer prefix inverse and the audited bound $N\le93$},
\path{reports/upper93_from94_20261016.tex}, 16 October 2026.
\bibitem{sharp} Sum-of-Subset research project,
\emph{A sharper non-progression lemma and the audited bound $N\le94$},
\path{reports/upper94_from95_20261016.tex}, 16 October 2026.
\bibitem{overlap} Sum-of-Subset research project,
\emph{Joint covering losses and the audited bound $N\le95$},
\path{reports/upper95_from96_20261016.tex}, 16 October 2026.
\bibitem{foundation} Sum-of-Subset research project,
\emph{The audited continuation from 115 to 103},
\path{reports/upper103_from115_20261014.tex}, 14 October 2026.
\bibitem{Lev1998} V.~F.~Lev, \emph{On consecutive subset sums},
Discrete Mathematics 187 (1998), 151--160.
\href{https://doi.org/10.1016/S0012-365X(98)80006-X}{doi:10.1016/S0012-365X(98)80006-X}.
\end{thebibliography}
\end{document}
